\section{An explicit quadratic relative simulation in iPRA}
\label{sec:quadratic-transfer}

Let $M_0$ denote the pinned primary PA194 table and let $M_L$ denote
PA194LC.  This section proves a relative result.  No assertion identifying
either machine's halting with a PA refutation is used as a premise.
For a literal table $M$, totalize its transition function by leaving a
halting configuration fixed, and write
\[
 H_M(t)\;:\!\iff\;
 \text{the configuration after $t$ transitions from the blank start is halted}.
\]
This is a decidable primitive-recursive predicate; in the totalized
semantics it means that the machine has halted by time $t$.

The exact startup and local-contract certificates for the two tables,
together with the symbolic-rule soundness argument below, give the
following ordinary computer-assisted arithmetic theorem:
\[
 \mathsf{iPRA}\vdash\forall t\,
 \bigl(H_{M_0}(t)\longrightarrow H_{M_L}(\Phi(t))\bigr),
 \qquad
 \mathsf{iPRA}\vdash\forall t\,
 \bigl(H_{M_L}(t)\longrightarrow H_{M_0}(\Phi(t))\bigr),
 \tag{Q}\label{eq:quadratic-transfer}
\]
where the explicit polynomial is
\[
 \boxed{\Phi(t)=58\,154+13(t+1)^2+681(t+1)
              =13t^2+707t+58\,848.}
\]
Consequently iPRA proves equivalence of their halting sentences, and
intuitionistic logic gives equivalence of their nonhalting sentences.
This is not a newly compiled Lean theorem.  The proof's mathematical
hypotheses are the complete literal-table contracts and their stated
arithmetical soundness, not the unfinished absolute consistency
calibration of the original primary table.

\subsection{Primitive-recursive configurations and local proofs}

Use the canonical tape zipper $(q,l,r)$, with the low bit of $r$ at the
head and the low bit of $l$ immediately to its left.  A write of $w$ followed
by a right move sends $(l,r)$ to
\[
 (2l+w,\lfloor r/2\rfloor),
\]
and a write followed by a left move sends it to
\[
 (\lfloor l/2\rfloor,
   4\lfloor r/2\rfloor+2w+(l\bmod 2)).
\]
This represents the entire finite-support tape, including arbitrary blank
exteriors, by natural numbers.  Table lookup, the totalized step, finite
word operations, and bounded iteration are primitive recursive.  Their
composition equations require only quantifier-free induction.

For each implementation augment the logical source configuration by a
nonnegative gap $g$, a first spare value $a_0$, and a finite spare tail
$A$.  The frame function $E(s)$ is thereby a primitive-recursive function,
not a relation requiring an unbounded existential witness.  The source
counter selects one of the finite list of valid instruction addresses;
the nine logical registers and the auxiliary fields are explicit coded
parameters.  Every frame has the nonhalting root control.  In particular,
a frame whose next source instruction is \emph{halt} is still a
nonhalting physical configuration before that instruction is executed.

The independent checker derives four generic list traversals directly
from the literal rows.  Their proofs use only: the defining equation for a
single transition; quantifier-free induction for a self-loop crossing a
constant-symbol run; finite coded-word identities; and induction on a
coded list's length.  If
\[
 B(v)=0\,1^{v+1},\qquad C(v)=1\,0\,1^v,
\]
the normalization identity $C(A)1=1B(A)$ follows from the corresponding
single-block identity by list-length induction.  The list traversal has
exact duration $\operatorname{width}(A)=\sum_{v\in A}(v+2)$.
Exterior-word codes remain free parameters throughout these proofs.
Canonical zipper equality replaces an unbounded assertion about every
tape coordinate.

The finite symbolic proof for each source site composes these lemmas.
Every decrement has its separate zero and successor-value branches.
All intermediate run lengths have coefficientwise nonnegative affine
expressions, and every complete contract has strictly positive duration.
The checker does not use arithmetic-box oracles.  Its finite successful
execution is evidence for the explicit derivations generated by these
sound rules; Python is not being treated as a proof-assistant kernel.

The two independent address resolvers agree with each other and with all
$3\,673$ entries, including jump targets, of the primary resolved source
program.  The implementations can have different physical paths, gaps,
and auxiliary updates.  Their logical projections are nevertheless the
same primitive-recursive source iteration, by quantifier-free induction.

\subsection{Synchronization and bounded extraction}

For one implementation, let $h(s)$ be the Boolean saying that the source
counter is at its explicit halt instruction.  Let $S$ be the augmented
source step, with a common absorbing logical halt extension and arbitrary
primitive-recursive auxiliary updates after such an instruction.  Write $s_n=S^n(s_0)$.  Let $k(s)\ge1$ be the exact
normal-instruction duration, and $k_H(s)\ge1$ the halt-instruction duration;
extend both functions totally on unused inputs.  If $b$ is the fixed
startup time, define
\[
 T(0)=b,\quad T(n+1)=T(n)+k(s_n),\qquad
 B(n)=\prod_{j<n}(1-h(s_j)).
\]
The Boolean $B(n)$ and clock $T(n)$ are primitive recursive.
The synchronization invariant is the quantifier-free implication
\[
 B(n)=1\ \longrightarrow\
 \bigl(\operatorname{Inv}(s_n)\ \wedge\
   \operatorname{run}_M(T(n),\mathrm{blank})=E(s_n)
   \ \wedge\ T(n)\ge b+n\bigr).
 \tag{S}\label{eq:source-synchronization}
\]
The startup is its base case.  The defining equation
$B(n+1)=B(n)(1-h(s_n))$ supplies both the predecessor invariant and the
nonhalt-source premise required by the next local contract.  Exact
step-iteration composition gives the physical equation.  Positivity of
$k$ gives the clock inequality.  The frame invariant is preserved because
the source targets explicitly preserve a first spare, keep the gap
nonnegative, and select another valid source address.  Thus the induction
formula contains no unbounded quantifier over runs, tapes, or proofs.

For the forward direction, a source witness $h(s_n)=1$ can be replaced by
the least $m\le n$ with that property, using bounded minimization.  Then
$B(m)=1$, and the halt-site contract gives a physical halt at
$T(m)+k_H(s_m)$.  This first-occurrence step avoids silently treating an
arbitrary source-halting witness as the first one.

For the reverse direction suppose $H_M(t)$.  If $B(t+1)=1$, equation
\eqref{eq:source-synchronization} gives a nonhalting root frame at time
\[
 T(t+1)\ge b+t+1>t.
\]
Absorbing halting implies that the machine is halted at this later time,
a contradiction.  Therefore $B(t+1)=0$.  A bounded first-hit function
returns some $m<t+1$ for which $h(s_m)=1$.  Its correctness is proved by
quantifier-free induction on the finite Boolean prefix.  In particular,
\[
 m\le t.
\]
There is no unbounded search, inverse decoding of intermediate tapes,
Markov principle, or elimination of double negation for an unbounded
existential.  The indexing convention matters: $h$ marks a
\emph{ready-to-execute} halt instruction, whose physical frame is still
nonhalting.  Under a different convention one may use the prefix through
index $t+1$ and the harmless bound $m\le t+1$ instead.

A hidden halt inside a claimed normal segment cannot defeat the
argument: in absorbing semantics a correct equation ending at a
nonhalting frame already excludes such an earlier halt.

\subsection{A uniform polynomial clock}

The actual complete startup tapes of both tables were separately
simulated and decoded.  After $b=58\,154$ transitions each is at native
control $84$, head $-2$, with source address $100000$ (instruction $25$),
and with
\[
 g=0,\qquad (r_0,\ldots,r_8)=(2,0,\ldots,0),
 \qquad a_0=1,\qquad A=[1,\ldots,1]\text{ of length }25.
\]
There are $26$ spare blocks altogether, so
$\operatorname{width}(A)=25(1+2)=75$.
The replay compares the entire finite support, not merely the control,
head position, or a bounded tape prefix.

For each of the $4\,780$ source outcomes of each literal table, a finite
coefficientwise comparison of its thirteen-slot duration vector gives
\[
 3\le k(s),k_H(s)\le
 534+8g+4\sum_{i=0}^{8}r_i+2a_0+2\operatorname{width}(A).
 \tag{D}\label{eq:affine-duration-envelope}
\]
Only the applicable normal or halt duration is used in a given case.
A successor-decrement certificate uses $r_i-1$ for its variable slot,
which can only decrease this upper bound.  The machine and contract
hashes are bound to the coefficient receipt.

Each source primitive increments at most one logical register by one.
Its gap update is $g'=g+\delta+o$, where both $\delta$ and the address
overflow bit $o$ belong to $\{0,1\}$.  Increment appends one spare value
$1$; every other primitive retains the spare list.  Consequently,
after $j$ normal instructions,
\[
 \sum_i r_i\le2+j,\qquad g\le2j,\qquad
 a_0=1,\qquad \operatorname{width}(A)\le75+3j.
\]
These are again quantifier-free induction invariants of the explicit
primitive target-frame updates.  Substitution in
\eqref{eq:affine-duration-envelope} gives
\[
 k(s_j),k_H(s_j)\le
 534+16j+4(2+j)+2+2(75+3j)=694+26j.
\]
Hence $n$ normal instructions take at most
\[
 b+\sum_{j<n}(694+26j)=b+13n^2+681n,
\]
and a first halt instruction at index $n$ yields halting by
\[
 b+13(n+1)^2+681(n+1).
\]
Given a physical halt at time $t$ in either implementation, bounded
extraction supplies the same logical halt-site index $n\le t$ for the
other.  Its forward bound and absorbing halting prove
\eqref{eq:quadratic-transfer}.  The affine upper bound is stronger than
needed for relative iPRA provability---total primitive-recursive clocks
and positivity already suffice---but makes both witness transformations
fully explicit.

The proof contains no compiler from arbitrary PA proofs to numerical
refutations.  That separate calibration is still needed to replace one
of the halting sentences by PA inconsistency.  Once it is established
for the primary table in iPRA, the relative theorem transfers it to
PA194LC immediately.
